% Generated draft: requires mathematical and source review.
\documentclass[11pt]{article}
\usepackage[margin=25mm]{geometry}
\usepackage{fontspec}
\setmainfont{Latin Modern Roman}
\usepackage{amsmath,amssymb,mathtools,unicode-math}
\setmathfont{Latin Modern Math}
\usepackage{enumitem,needspace,xurl,hyperref}
\setlength{\parindent}{0pt}
\setlength{\parskip}{5pt}
\setlength{\emergencystretch}{3em}
\newcommand{\N}{\mathbb{N}}
\newcommand{\Z}{\mathbb{Z}}
\newcommand{\Q}{\mathbb{Q}}
\newcommand{\R}{\mathbb{R}}
\newcommand{\C}{\mathbb{C}}
\newcommand{\F}{\mathbb{F}}
\newcommand{\A}{\mathbb{A}}
\newcommand{\PP}{\mathbb{P}}
\newcommand{\E}{\mathbb{E}}
\newcommand{\eps}{\varepsilon}
\newcommand{\dd}{\,\mathrm{d}}
\begin{document}
\section*{UnsolvedMath ID: 30002583}
\textbf{Source problem code:} catalogue entry
\subsection{Definitions and mathematical statement}
Let \(M\) be a \(W^*\)-algebra and let \(E\) be a \(W^*\)-correspondence over \(M\), that is, a self-dual right Hilbert \(M\)-module with a normal left action of \(M\) by adjointable operators. Let \(Z\) be an admissible sequence of weights defining the weighted tensor algebra \(\mathcal T_+(E,Z)\), as intended in the supplied record. The record does not specify the admissibility conditions or the construction of this algebra.

Write \(\Sigma\) for the intended index class of representations \(\sigma:M\to B(H_\sigma)\), and write \(D_{X,\sigma}\) for their associated weighted intertwiner domains, retaining the notation of the original statement. The representation class, the domains, and the relation between the parameter \(X\) and the weights \(Z\) are not defined in the supplied record. For \(F\in\mathcal T_+(E,Z)\), denote its asserted Berezin transform by
\[
\widehat F=(\widehat F_\sigma)_{\sigma\in\Sigma},\qquad
\widehat F_\sigma:D_{X,\sigma}\longrightarrow B(H_\sigma).
\]
Let \(\mathfrak M(D)\) denote the intended class of matricial families \(f=(f_\sigma)_{\sigma\in\Sigma}\) on these domains. This notation is a placeholder for the source's unspecified compatibility and regularity requirements, not a definition of those requirements.

The supplied premise is that \(\widehat F\in\mathfrak M(D)\) for every \(F\in\mathcal T_+(E,Z)\). The requested converse is the following realization question: for every choice of \(M,E,Z\) satisfying the intended hypotheses, is it true that
\[
\forall f\in\mathfrak M(D)\quad
\exists F\in\mathcal T_+(E,Z)\quad
\forall\sigma\in\Sigma\quad
\forall\mathfrak z\in D_{X,\sigma},\qquad
f_\sigma(\mathfrak z)=\widehat F_\sigma(\mathfrak z)?
\]
Equivalently, is the asserted Berezin-transform map \(\mathcal T_+(E,Z)\to\mathfrak M(D)\) surjective? No uniqueness is requested. A fully specified mathematical assertion requires the missing definitions of admissibility, the representation class, the weighted domains, and the intended matricial-family conditions. In particular, the record does not state whether boundedness, continuity, analyticity, or boundary-extension conditions are imposed on the families.
\subsection{Short English statement}
The question asks whether every intended matricial family on the weighted intertwiner domains is the Berezin transform of a single element of \(\mathcal T_+(E,Z)\). The supplied record leaves essential domain and family conditions unspecified, so the displayed realization question remains conditional on those definitions. Its supplied \texttt{partially\_solved} label and unweighted comparison are dated curation requiring independent review.
\subsection{Sources}
\begin{itemize}
\item Supplied local UnsolvedMath registry record, permanent ID \(30002583\), title \emph{Realization of Matricial Families by Tensor-Algebra Elements}, exact source problem code \texttt{catalogue entry}, in the collection \emph{Oberwolfach Reports Open Problems}. The record gives proposed year 2014 and names no proposer.
\item Supplied report citation: \emph{Hilbert Modules and Complex Geometry} (2015), Oberwolfach Reports, pp.\ 1139--1219, DOI \texttt{10.4171/owr/2014/21}. \url{https://doi.org/10.4171/owr/2014/21}. The background describes these reports as records of mathematical workshops and their open-problem discussions.
\item Supplied literature reference: \emph{Tensorial function theory}. \url{https://arxiv.org/abs/1210.2964}.
\item Supplied literature reference: \emph{Published tensorial function theory}. \url{https://doi.org/10.1007/s00020-013-2062-4}.
\item Supplied literature reference: \emph{Oberwolfach report}. \url{https://doi.org/10.4171/owr/2014/21}.
\item Literature assessment dated exactly \texttt{2026-08-21}: the supplied curation labels the problem partially solved and states that the converse characterization is proved in the unweighted case, whereas the weighted-correspondence version is not settled by the cited theory. This is supplied dated curation derived from the supplied OWR material, not an independently verified assessment; its time-sensitive status and cited claims require independent verification.
\item The registry labels the statement \texttt{corrected\_verified} and says that mathematical content was preserved while notation, typography, or omitted definitions were repaired for publication. These are supplied editorial claims, not independent mathematical verification of this draft or of the cited sources.
\end{itemize}
\end{document}
